\chapter{Correlated electrons}
\setwlchapter{correlated-electrons}

What does it cost to move one electron between two hydrogen-like atoms?
At large separation and with one electron on each atom, the final state
contains an empty orbital and a doubly occupied one. Creating this pair
costs the local repulsion $U$. Yet a virtual excursion into precisely the
same state lowers the energy of a spin singlet. A charge excitation at
energy $U$ and a spin excitation at energy of order $t^2/U$ thus have
a common origin. The two-site Hubbard model lets us follow both without
making an approximation. We will use its exact spectrum to judge two
quite different approximations: eliminating charge fluctuations, and
replacing the interaction by a self-consistent field.

The examples form one session. Start a fresh kernel with the SNEG paclet
installed and evaluate the inputs in order. The convention for a level
is \fn{c[i]}; a single electron operator carries two additional indices,
as in \fn{c[CR,i,UP]}. Spin projections are in units of $\hbar$, and
$k_{\mathrm B}=1$.

\section{One orbital: charge, spin and particle--hole symmetry}
For one Wannier orbital on each of $L$ sites we use
\begin{equation}
 H=-t\sum_{\langle ij\rangle,\sigma}
 (c^\dagger_{i\sigma}c_{j\sigma}+c^\dagger_{j\sigma}c_{i\sigma})
 +U\sum_i D_i,\qquad D_i=n_{i\uparrow}n_{i\downarrow}.
 \label{eq:hubbard}
\end{equation}
Here $n_{i\sigma}=c^\dagger_{i\sigma}c_{i\sigma}$,
$\sigma=\uparrow,\downarrow$, and a bond occurs once in the sum.
Throughout this chapter $t$ is real and $U>0$. In a crystal $U$ is the
interaction screened by degrees of freedom excluded from this orbital
model; it need not equal an isolated-atom Coulomb integral.
Our starting point is the atomic-limit construction in
Ref.~\cite{hubbard}.

\begin{wl}
Needs["sneg`"];
snegfermionoperators[c, d];
snegrealconstants[t, U, \[Epsilon]];
atom = \[Epsilon] number[d[]] + U hubbard[d[]];
{number[d[]], hubbard[d[]]}
\end{wl}

The minus sign that may appear when $D_i$ is printed in normal order is
not an attractive interaction. Moving an annihilator through a creator
of the other spin costs one interchange. A sign-independent way to read
the result is to act on states. The following basis is labelled by
$Q=N-1$ and total spin $S$. Each sector of this atom contains a single
multiplet; its energy is the sole entry of the block.
\begin{wl}
ba = qsbasisvc[{d[]}];
rows = {Sequence @@ First[#], #[[2, 1, 1]]} & /@ makematricesbzvc[atom, ba];
TableForm[rows, TableHeadings -> {None, {"Q", "S", "Energy"}}]
\end{wl}

The empty state has energy zero, the two spin states have energy
$\epsilon$, and $\ket{2}=d^\dagger_\uparrow d^\dagger_\downarrow\ket{0}$
has energy $2\epsilon+U$. \fn{qsbasisvc} stores one representative of
each spin multiplet: the dimension of a block counts copies of the
representation, not its $2S+1$ magnetic substates.
With $\epsilon$ measured relative to the chemical potential, the atom
has a spin-doublet ground state for $-U<\epsilon<0$. At either
endpoint two charge sectors cross; differentiating a single eigenvalue
through such a crossing would give the wrong ground-state occupation.

At half filling a particle--hole transformation on a bipartite lattice
takes $c_{i\sigma}\mapsto\eta_i c^\dagger_{i\sigma}$, with $\eta_i=+1$
on one sublattice and $-1$ on the other. Hopping is invariant because
neighbours have opposite $\eta$ and distinct fermions anticommute.
Locally $n_i\mapsto2-n_i$ and $D_i\mapsto1-n_i+D_i$. Thus the
grand Hamiltonian with $\epsilon=-U/2$ is invariant. Equivalently,
\[
 U(n_{i\uparrow}-\tfrac12)(n_{i\downarrow}-\tfrac12)
 =U D_i-\tfrac U2 n_i+\tfrac U4 .
\]
Dropping the last term changes the energy origin; dropping the term
linear in $n_i$ also changes the chemical potential. Neither affects
splittings at fixed $N$, but only a constant can be ignored when
comparing absolute addition and removal energies at a specified $\mu$.

\subsection{A free spin without a static polarization}
At $\epsilon=-U/2$ the empty and doubly occupied states both have
energy zero and the two singly occupied states each have energy $-U/2$.
For $T>0$ the grand partition function is $2(1+e^{U/(2T)})$. It
follows that the mean occupation is exactly one at \emph{any}
temperature; it tells us nothing about moment formation. A more
discriminating observable is $4\langle(S^z)^2\rangle$, which equals
the probability of single occupation. The complementary observable,
$\langle(n-1)^2\rangle$, counts empty and doubly occupied sites.
Test the two in a thermal trace over all four charge and spin states:
\begin{wl}
sz2 = pow[number[d[], UP] - number[d[], DO], 2];
q2 = pow[number[d[]] - 1, 2];
Expand[sz2 + q2]
\end{wl}
\begin{wl}
hat = atom /. {\[Epsilon] -> -1, U -> 2};
rows = Table[Prepend[Round[thermalaverage[#, hat, {d}, temp], 1/1000] & /@
    {number[d[]], sz2, q2}, temp], {temp, {0, 1, 4}}];
TableForm[N[rows], TableHeadings -> {None, {"T", "Charge", "Moment", "Charge variance"}}]
\end{wl}
The columns are $T$, $\langle n\rangle$, $4\langle(S^z)^2\rangle$ and
$\langle(n-1)^2\rangle$, for $U=2$ in the energy units used here.
At zero temperature \fn{thermalaverage} averages equally over the two
ground states, so $\langle S^z\rangle=0$ although the spin moment is
maximal. A field term $-h S^z$ lifts that degeneracy. Since $S^z$
commutes with the atomic Hamiltonian, the static susceptibility at
$h=0$ is $\chi=\langle(S^z)^2\rangle/T$ (the magnetic moment multiplying
the field is set to one). At low nonzero $T$ the free-spin contribution
$1/(4T)$ survives; at high $T$ the two spinless states halve its
coefficient. Taking $T\to0$ at fixed $h>0$ selects a pure spin state
with a different magnetization from the zero-field mixture.

\section{Exact solution of the Hubbard dimer}
Connect two sites by \emph{one} bond. The empty and doubly occupied
local states have spin zero, so each can mix with a two-site singlet
but not with a triplet. Build the full Hamiltonian before restricting
the basis. Products of operators are made with \fn{nc}; ordinary
multiplication suffices for a scalar times an operator.
This construction recasts the two-site symmetry problem of Steeb and
Hardy~\cite[Sec.~3.4, problem 84]{steeb} in a basis of total spin
instead of fixed spin projection.
\begin{wl}
sites = {c[1], c[2]};
h0 = U Total[hubbard[sites]];
v = -t hop[c[1], c[2]];
hd = h0 + v;
Expand[{komutator[hd, Total[number /@ sites]],
  komutator[hd, Total[spinplus /@ sites]]}]
\end{wl}
The second zero checks more than conservation of $S^z$: together with
Hermiticity it establishes invariance under spin rotations. At $L=2$,
$Q=N-2$, so the half-filled blocks are those with $Q=0$.
\begin{wl}
bd = qsbasisvc[sites];
sector[s_] := Select[bd, First[#] == {0, s} &];
mat[h_, b_] := First[makematricesbzvc[h, b]][[2]];
hs = mat[hd, sector[0]];
{Dimensions[hs], mat[hd, sector[1]]}
\end{wl}
\begin{wl}
hs // MatrixForm
\end{wl}
In this phase convention the diagonal states, in order, are a doublon
on one site, the spin singlet with one electron on each site, and a
doublon on the other site. Each amplitude $-\sqrt{2}\,t$ is the sum of
two indistinguishable hopping paths out of the normalized singlet.
Rephasing one basis vector flips a row and the corresponding column;
it cannot change an eigenvalue. With $P(x)=\det(H_{S=0}-xI)$:
\begin{wl}
Factor[CharacteristicPolynomial[hs, x]]
\end{wl}
The antisymmetric doublon combination decouples at energy $U$. The
symmetric combination couples to the singly occupied singlet by $-2t$.
The other two singlet energies and the singlet--triplet splitting are
\begin{equation}
 E_\pm=\frac{U\pm\sqrt{U^2+16t^2}}2,\qquad
 \Delta_s=E_T-E_-=-E_- .
 \label{eq:dimer}
\end{equation}
At $U=0$ the singlet fills the bonding orbital with both spins and has
energy $-2|t|$. At $t=0$ the singly occupied singlet and the triplet
are degenerate. Turning on hopping at large $U$ leaves the triplet at
zero but lowers the singlet. This is the simplest example of kinetic
energy selecting antiferromagnetic correlations.

\subsection{What remains of double occupation?}
The derivative $\partial E_-/\partial U=\langle D_1+D_2\rangle$
counts the doublons in the exact ground state. The equality follows by
differentiating the eigenvalue equation for a normalized eigenvector;
terms involving its derivative cancel. It is valid away from a crossing.
\begin{wl}
emin = (U - Sqrt[U^2 + 16 t^2])/2;
dsite = Simplify[D[emin, U]/2];
{dsite, dsite /. {U -> 0, t -> 1}}
\end{wl}
For $U\gg |t|$ the probability per site is $2t^2/U^2$ to leading order.
It is small but cannot be set to zero before evaluating the energy.
In fact $\langle v\rangle=t\,\partial E_-/\partial t$, and at this order
the kinetic gain is twice the interaction cost. Setting $U=4t$ gives
an instructive intermediate case: the exact singlet still has no local
magnetization, even though its local moment
$\langle(n_{i\uparrow}-n_{i\downarrow})^2\rangle=1-2\langle D_i\rangle$
already substantially exceeds the free-electron value $1/2$.

\begin{pitfall}
Occupation vectors refer to the current list of modes. A new call to a
basis generator can change it. Construct the matrix before creating a
basis on different sites, or regenerate the original basis first.
The numerical matrix, once constructed, is independent of that list.
\end{pitfall}

\section{Eliminating virtual charge excitations}
Let $P$ select one electron per site and $\bar P=1-P$. In the half-filled
dimer $P H_0 P=P v P=0$. Every intermediate state connected to $P$ by
one hop has energy $U$, so to second order
\[
 H^{(2)}=-P v\bar P vP/U .
\]
It is negative semidefinite. In a spin triplet there is no accessible
doublon, hence no shift. For a singlet the squared norm of $\bar P v
\ket{s}$ is $4t^2$. Since the singlet projector is
$1/4-\vS_1\cdot\vS_2$, this already fixes the entire spin Hamiltonian.

\begin{wl}
p = nc @@ (projector1 /@ sites);
heff = schriefferwolff[h0, v, p, 2];
snegxoperators[X, "HubbardAtom"];
snegspinoperators[{S, 1/2}];
Expand[xtospin[fermionstox[heff[[3]], c -> X], X -> S]]
\end{wl}

Read the transverse terms with
$S_1^xS_2^x+S_1^yS_2^y=(S_1^+S_2^-+S_1^-S_2^+)/2$.
The result is $J(\vS_1\cdot\vS_2-1/4)$ with $J=4t^2/U$.
\fn{heff[[3]]} is order two: the returned list starts with
order zero. The constant $-J/4$ matters when comparing energies
rather than splittings. A controlled expansion requires $|t|/U\ll1$.
The exact dimer shows both the sign and scale of the first omitted term:
\begin{wl}
Normal[Series[-emin, {t, 0, 4}, Assumptions -> U > 0]] //
  Simplify[#, U > 0] &
\end{wl}
Figure~\ref{fig:dimer-exchange} compares the two truncations. The
fourth-order correction lowers the second-order splitting. At $U=2t$
it cancels the entire leading term, although the exact gap is positive.
The relevant scale is visible already in Eq.~\eqref{eq:dimer}: the
Taylor series of the square root about $t=0$ converges for
$|4t/U|<1$, not throughout the range in the figure.

\begin{figure}[htbp]
\centering
\begin{tikzpicture}
\begin{axis}[width=.95\linewidth,height=53mm,xmin=2,xmax=12,ymin=0,ymax=2.05,
  xlabel={$U/t$},ylabel={$\Delta_s/t$},tick label style={font=\small},
  legend cell align=left,legend style={font=\small,draw=none,fill=none},
  samples=100,domain=2:12]
\addplot[very thick,snegblue] {8/(sqrt(x*x+16)+x)};
\addlegendentry{exact}
\addplot[black,dashed,thick] {4/x};
\addlegendentry{second order}
\addplot[black,densely dotted,thick] {4/x-16/(x*x*x)};
\addlegendentry{fourth order}
\end{axis}
\end{tikzpicture}
\caption{Singlet--triplet gap of a half-filled dimer, $t>0$. All
curves use the same one-bond hopping. At large $U/t$ the remaining
error after the fourth-order correction is of order $(t/U)^5$ in
the plotted units.}\label{fig:dimer-exchange}
\end{figure}

\subsection{Doping: when the first-order term survives}
Remove an electron. The low-energy subspace must now include holes, so
use $P_G=\prod_i(1-D_i)$ rather than the projector onto exactly one
electron per site. A hop onto an empty neighbour need not leave $P_G$.
With $\tilde c_{i\sigma}=c_{i\sigma}(1-n_{i\bar\sigma})$, the
usual nearest-neighbour $t$--$J$ Hamiltonian is
\begin{align}
 H_{tJ}={}&-t\sum_{\langle ij\rangle\sigma}
  (\tilde c^\dagger_{i\sigma}\tilde c_{j\sigma}+\hc)\nonumber\\
 &+J\sum_{\langle ij\rangle}
   (\vS_i\cdot\vS_j-\tfrac14 n_i n_j).
\end{align}
The density factor makes the exchange vanish if either site is empty.
On three or more sites this expression does \emph{not} exhaust
order $t^2/U$. A hop $k\to j$ can create a doublon, followed by a hop
$j\to i$ that removes it with $i\ne k$. If $i$ was initially empty,
the final state is in $P_G$, but the hole has moved from $i$ to $k$.
For real nearest-neighbour hopping the extra term can be written
\[
 -\frac{t^2}{U}\!
 \sum_{j,\;i\ne k\in\mathcal{N}(j),\;\sigma}
 P_G\big(c^\dagger_{i\sigma}n_{j\bar\sigma}c_{k\sigma}
 -c^\dagger_{i\sigma}c^\dagger_{j\bar\sigma}
 c_{j\sigma}c_{k\bar\sigma}\big)P_G .
\]
The neighbours $(i,k)$ are ordered; the reverse path is already in
the sum. The minus sign inside the parentheses distinguishes the two
possible spin transfers. At half filling both transfers to a third site
are blocked by the final projector. Away from half filling, neglecting
them is an additional modelling decision, not a higher-order truncation.

\paragraph{Isolating the three-site matrix element.}
Take an open chain $1$--$2$--$3$ with two electrons in a spin singlet.
There are three states with the electrons on different sites and three
with both electrons on one site. A singlet shared by sites 1 and 2 can
turn into a singlet on 2 and 3 in two hops through a doublon on site 2.
In contrast, exchange on either individual bond cannot change the hole
position. This makes the off-diagonal element of $H^{(2)}$ between
the two bond singlets a diagnostic for the three-site term.
\begin{wl}
line3 = c /@ Range[3];
b3 = Select[qsbasisvc[line3], First[#] == {-1, 0} &];
nd = Diagonal[mat[Total[hubbard[line3]], b3]];
lo = Flatten[Position[nd, 0]]; hi = Flatten[Position[nd, 1]];
v3 = mat[-hop[c[1], c[2]] - hop[c[2], c[3]], b3];
k3 = -v3[[lo, hi]] . v3[[hi, lo]];
k3 // MatrixForm
\end{wl}
The printed matrix is in units of $t^2/U$; the dimensionless entries
are exact. A pair occupying sites 1 and 3 has no virtual doublon one
hop away, hence a zero row. The two adjacent singlets have diagonal
shift $-4$ and mutual coupling of magnitude $2$: only the intermediate
doublon on the shared site connects them. Which row is zero and the
sign of the coupling depend on the chosen basis. The eigenvalues do
not. Removing the three-site term from this matrix gives a
spurious degeneracy:
\begin{wl}
{Sort[Eigenvalues[k3]], Sort[Diagonal[k3]]}
\end{wl}
These numbers diagnose the second-order \emph{operator}, not the full
energy shifts of the doped chain. To obtain the low-energy spectrum
through this order, add the first-order block $t\,v_{3,\mathrm{lo},\mathrm{lo}}$
to $(t^2/U)\,k_3$ and diagonalize. At half filling the first-order
block would be zero.

\subsection{Four sites: a process with no dimer counterpart}
Four successive hops around a square can restore single occupation
while permuting all four spins. On a bipartite plaquette, returning to
single occupation requires an even number of hops, so the first
corrections to superexchange occur at order $t^4/U^3$. They include
both pair exchanges and a cyclic permutation plus its inverse
\cite{macdonald}. Renormalizing $J$ alone cannot describe an entire
plaquette spectrum. A coefficient quoted for a four-spin product is
not directly the coefficient of a permutation: expanding a spin-$1/2$
swap, $P_{ij}=2\vS_i\cdot\vS_j+1/2$, generates lower-body terms too.

The plaquette example compares the lowest sixteen half-filled energies
with the sixteen energies in the projected spin space. From the
checkout run:
\begin{lstlisting}[style=wl]
wolframscript -file examples/schriefferwolff-plaquette.m
\end{lstlisting}
Keeping $t$ fixed and doubling large $U$ reduces the leading
absolute spectral error by about eight at second order ($t^4/U^3$)
and by about thirty-two at fourth order ($t^6/U^5$). Matching the
dimension and the energy origin is essential in this comparison.

\begin{furtherreading}
For the passage from electrons to spins see Auerbach~\cite{auerbach}.
In a junction the virtual states can include superconducting
quasiparticles; their energies replace the single denominator $U$ of
the dimer. An application is the Josephson harmonic expansion of
Bácsi, Iličin and Žitko~\cite{bacsi}.
\end{furtherreading}

\section{Measuring a charge gap on a finite cluster}
Define $E_0(N)$ as the lowest energy among \emph{all} spin sectors
at particle number $N$. At zero temperature that charge sector is stable
over a chemical-potential interval of width
\begin{equation}
 \Delta_c=\mu_+-\mu_-
 = E_0(N+1)+E_0(N-1)-2E_0(N).
 \label{eq:chargegap}
\end{equation}
This second difference eliminates every term linear in $N$. It is
not the energy of the lowest neutral excitation: the dimer, for
instance, admits an arbitrarily soft spin excitation as $t^2/U\to0$,
even as its charge gap grows.

Take four sites on a ring, count each bond once, and measure energies in
units of $t>0$. Here we keep $S^z$, rather than total $S$, and use all
magnetic substates. The blocks are larger but their dimensions now sum
directly to the Fock-space dimension $4^4$.
\begin{wl}
ring = c /@ Range[4];
hr = -Sum[hop[c[i], c[Mod[i, 4] + 1]], {i, 4}] +
  U Total[hubbard[ring]];
br = qszbasisvc[ring];
blocks = makematricesbzvc[hr, br];
{Length[blocks], Total[Length /@ blocks[[All, 2]]]}
\end{wl}
\begin{wl}
e0[q_, u_] := Min[Flatten[Eigenvalues[N[#[[2]] /. U -> u]] & /@
  Select[blocks, #[[1, 1]] == q &]]];
gap[u_] := e0[1, u] + e0[-1, u] - 2 e0[0, u];
rows = Table[{u, N[Round[gap[u], 1/1000]]}, {u, {0, 2, 4, 8, 16}}];
TableForm[rows, TableHeadings -> {None, {"U/t", "Charge gap / t"}}]
\end{wl}
Each row gives $(U/t,\Delta_c/t)$ rounded to three decimal places.
At $U=0$ the allowed momenta are $0,\pm\pi/2,\pi$ (lattice spacing
one). The two zero-energy orbitals make this an open shell, so adding
and removing an electron have identical one-particle thresholds.
The same calculation on a closed-shell cluster can yield a nonzero
gap in a metal. Inferring a bulk gap requires a sequence of sizes
and control of such shell effects.

For comparison, at $U=0$ alternate the hopping amplitudes
$t_1,t_2>0$ in a chain with two sites per cell. The Bloch energies are
$\pm|t_1+t_2 e^{ik}|$, where $k$ is in inverse cell lengths. The
half-filled band gap $2|t_1-t_2|$ persists in the thermodynamic limit,
without an interaction. In the large-$U$ uniform Hubbard chain charge
motion instead costs a doublon energy even in the absence of a
one-particle band gap. The charge-gap measurement alone cannot tell
these mechanisms apart; compare the noninteracting spectrum, double
occupancy and spin response.

\subsection{A shell-effect check without many-body diagonalization}
The noninteracting uniform ring is a useful calibration. Sort its
$L$ one-particle levels \emph{per spin}; at half filling each spin
occupies $L/2$ levels. The finite-size addition gap is the difference
of two consecutive levels at the Fermi surface.
\begin{wl}
shell[l_] := With[{band = SortBy[Table[-2 Cos[2 Pi j/l], {j, 0, l-1}], N]},
  band[[l/2 + 1]] - band[[l/2]]];
TableForm[Table[{l, N[shell[l]]}, {l, {4, 6, 8, 10}}],
  TableHeadings -> {None, {"Sites", "Charge gap / t"}}]
\end{wl}
For $L=4p$, both Fermi points lie on the grid and the gap is zero.
For $L=4p+2$, the grid misses each Fermi point by $\pi/L$, so
$\Delta_c/t=4\sin(\pi/L)$: a finite gap at \emph{every} size on this
subsequence tends to zero as $1/L$. In an interacting calculation one
cannot simply subtract this formula; the velocities and the many-body
levels have changed. But treating the two subsequences separately
prevents an alternation of boundary-condition effects from masquerading
as an interaction threshold.

\section{The atomic propagator and Hubbard-I}
Consider the retarded one-electron Green's function
$G_\sigma(z)=\langle\!\langle d_\sigma;d^\dagger_\sigma
\rangle\!\rangle_z$, where $z=\omega+i0^+$. Its spectral density
$A_\sigma(\omega)=-\operatorname{Im}G_\sigma(\omega+i0^+)/\pi$
has unit integral \emph{per spin}. Use the same grand Hamiltonian
as in the atomic spectrum. Since $n_{\bar\sigma}$ commutes with it,
the two operators $d_\sigma(1-n_{\bar\sigma})$ and
$d_\sigma n_{\bar\sigma}$ evolve at fixed frequencies $\epsilon$ and
$\epsilon+U$, respectively. The equation-of-motion hierarchy closes:
\begin{align*}
 (z-\epsilon)G_\sigma &= 1+U F_\sigma,\\
 (z-\epsilon-U)F_\sigma &= \langle n_{\bar\sigma}\rangle ,
\end{align*}
where $F_\sigma=\langle\!\langle d_\sigma n_{\bar\sigma};
d^\dagger_\sigma\rangle\!\rangle_z$. These are operator identities;
no factorization of an average was used.
\begin{wl}
gatom = greensfunction[d[AN, UP], d[CR, UP], atom, z];
g = gatom /. avg[number[d[], DO]] -> w;
Apart[g, z]
\end{wl}
Writing $w=\langle n_{\downarrow}\rangle$, the poles and their weights
are transparent in
\begin{equation}
 G_\uparrow(z)=\frac{1-w}{z-\epsilon}
              +\frac{w}{z-\epsilon-U}.
 \label{eq:atomicG}
\end{equation}
The upper pole exists only when a down electron is available on the
site. With $\epsilon=-U/2$ and a spin-symmetric ensemble, $w=1/2$;
the spectral peaks then have equal weight and are separated by $U$.
A pure member of the atomic doublet has the same energy but different
spin-resolved weights. Substituting $w=1/2$ is a specification of the
state, not a step in the equation-of-motion calculation.
\begin{wl}
ph = Factor[g /. {\[Epsilon] -> -U/2, w -> 1/2}];
{ph, Limit[z ph, z -> Infinity]}
\end{wl}

\paragraph{Checking both residues against a trace.}
The half-filled atom is a poor test of thermal weights: symmetry
forces them both to $1/2$ regardless of temperature. Instead take
$\epsilon=-1/4$, $U=2$, and $T=1/2$. Removing or adding one spin-up
electron still gives only the two transition frequencies $\epsilon$
and $\epsilon+U$, but their residues differ. The Lehmann representation
sums matrix elements between exact many-body eigenstates. Each row
below is $(E_n-E_m,w_{mn})$ after collecting equal frequencies:
\begin{wl}
asym = atom /. {\[Epsilon] -> -1/4, U -> 2};
poles = lehmann[d[AN, UP], d[CR, UP], asym, {d}, 1/2];
TableForm[N[Round[poles, 1/1000]], TableHeadings -> {None, {"Frequency", "Residue"}}]
\end{wl}
For a retarded anticommutator the factor in the residue is the
\emph{sum} of the two Boltzmann probabilities, not their difference.
The atomic partition function is
$Z=1+2e^{-\epsilon/T}+e^{-(2\epsilon+U)/T}$.
The upper peak has weight $w=(e^{-\epsilon/T}+e^{-(2\epsilon+U)/T})/Z$.
These facts give two independent checks: the total weight is one,
and the upper weight equals the mean down-spin occupation. Compare the
full complex propagators at a frequency off the real axis, where a
pole neither diverges nor needs artificial broadening:
\begin{wl}
wth = thermalaverage[number[d[], DO], asym, {d}, 1/2];
z0 = 1/3 + I/4;
gtrace = Total[poles[[All, 2]]/(z0 - poles[[All, 1]])];
{Abs[1 - Total[poles[[All, 2]]]] < 10^-12,
 Abs[gtrace - (g /. {z -> z0, \[Epsilon] -> -1/4, U -> 2, w -> wth})] < 10^-12}
\end{wl}

\subsection{What changes when the atom is put back in the lattice?}
Define the local self-energy by
$G_\mathrm{at}^{-1}(z)=z-\epsilon-\Sigma_\mathrm{at}(z)$. Inverting
Eq.~\eqref{eq:atomicG} gives
\[
 \Sigma_\mathrm{at}(z)=Uw+
 \frac{U^2 w(1-w)}{z-\epsilon-U(1-w)} .
\]
The static term is the Hartree shift. The pole in the second term is
what opens a correlation gap in this approximation. Hubbard-I inserts
this atomic self-energy into the lattice propagator, or equivalently
uses $G_k(z)=[G_\mathrm{at}^{-1}(z)-\varepsilon_k]^{-1}$, with
$\varepsilon_k$ the dispersion of the hopping term alone.
\begin{wl}
gk = Together[ph/(1 - ek ph)];
Factor[Denominator[gk]]
\end{wl}
At particle--hole symmetry its poles are
$z_\pm=(\varepsilon_k\pm\sqrt{\varepsilon_k^2+U^2})/2$.
For $U>0$ they lie on opposite sides of zero. For $U=0$ the apparent
additional pole cancels against a zero in the numerator, leaving
$1/(z-\varepsilon_k)$. Checking residues as well as pole locations
prevents a spurious flat band. The construction is exact in the atomic
limit and recovers free propagation at $U=0$. It omits intersite
spin correlations and has no coherent Fermi-liquid resonance. Its
half-filled gap at arbitrarily small $U>0$ does not establish a Mott
transition for a general lattice.

\section{Can a static field account for the singlet energy?}
In a collinear Hartree state, write $n_{i\sigma}=\bar n_{i\sigma}
+\delta n_{i\sigma}$ and discard
$\delta n_{i\uparrow}\delta n_{i\downarrow}$. The remaining interaction is
\begin{equation}
 U\sum_i(\bar n_{i\downarrow} n_{i\uparrow}+
          \bar n_{i\uparrow} n_{i\downarrow}-
          \bar n_{i\uparrow}\bar n_{i\downarrow}).
\end{equation}
The last term does not change the orbitals, but omitting it counts
the interaction energy twice. In a translationally invariant metal,
a small uniform polarization $m=\bar n_\uparrow-\bar n_\downarrow$
splits the two spin bands by $Um$. For a smooth one-spin density of
states $\rho(E_F)$ per site, the induced polarization at zero
temperature is $\rho(E_F)Um$ to first order. Thus the unpolarized
solution is unstable to \emph{uniform} magnetization for
$U\rho(E_F)>1$ (Stoner criterion). It need not be the first
instability: the susceptibility at nonzero wave vector can be larger.

On the dimer, try staggered polarization instead:
$m_1=-m_2=m$, with one electron per site on average. Apart from the
common Hartree shift, each spin sees diagonal potentials $\pm Um/2$
and off-diagonal hopping $-t$. Filling the lower eigenstate for each
spin gives
\[
 m=\frac{Um/2}{\sqrt{t^2+(Um/2)^2}} .
\]
Besides $m=0$, for $U>2|t|$ there are two solutions related by spin
reversal with $m=\pm\sqrt{1-(2t/U)^2}$. Below we seed one branch;
the solver adjusts the chemical potential for two particles.
\begin{wl}
hf = meanfield[hd /. {t -> 1, U -> 4}, "Hartree"];
seed = {avg[number[c[1], UP]] -> 0.9, avg[number[c[2], DO]] -> 0.9};
sol = meanfieldsolve[hf, sites, seed, "ParticleNumber" -> 2];
m1 = avg[number[c[1], UP]] - avg[number[c[1], DO]];
{sol["Converged"], Round[sol["Energy"], 1/1000],
  Round[m1 /. sol["Averages"], 1/1000] // N}
\end{wl}
At these parameters the unpolarized determinant has energy
$-2t+U/2=0$, the staggered determinant has energy $-t/2$, and the
exact energy is $(2-2\sqrt{2})t$. Improving the variational energy by
allowing a staggered field has introduced a false broken symmetry on a
finite dimer. In particular $m_1$ is not the local fluctuating moment
calculated in the singlet. At large $U$, Hartree theory yields
$-2t^2/U$ for this determinant; the exact singlet gains twice as much
through transverse spin exchange. Convergence of the iteration is not
a test of either the symmetry or the accuracy of the state.

\section{Local constraints: Hubbard operators and auxiliary particles}
Operator projection is easiest to understand on one site. For
$a,b\in\{0,\uparrow,\downarrow,2\}$ let $X^{ab}=\ket{a}\bra{b}$.
Then $X^{ab}X^{cd}=\delta_{bc}X^{ad}$ and $\sum_a X^{aa}=1$.
The creation operator $d^\dagger_\uparrow$ is the sum of the transitions
$X^{\uparrow0}+X^{2\downarrow}$. For down spin the second transition
has a minus sign because $\ket{2}$ was defined with the up creator
on the left. Work with the corresponding fermionic polynomials to
test the projected algebra:
\begin{wl}
xi = xoperator[d[], "0", "UP"];
Expand[antikomutator[xi, conj[xi]]]
\end{wl}
The anticommutator is $X^{00}+X^{\uparrow\uparrow}
=1-n_\downarrow$, not the identity. The missing weight is precisely
the transition at the upper Hubbard pole. At distinct sites, odd
transitions anticommute: ordinary commuting site matrices would lose
the fermion sign.
\begin{wl}
pg = 1 - hubbard[d[]];
Expand[{nc[pg, pg] - pg, nc[pg, d[AN, UP], pg] - xi}]
\end{wl}

One can instead keep canonical auxiliary fields at the price of a
constraint. In the no-doublon space introduce a boson $b$ for
$\ket0$ and fermions $f_\sigma$ for $\ket\sigma$. Identify
$X^{0\sigma}=b^\dagger f_\sigma$ and impose
\[
 b^\dagger b+\sum_\sigma f^\dagger_\sigma f_\sigma=1
 \quad\text{on every site}.
\]
The three physical states are precisely the states with one auxiliary
particle. The substitutions $b\to e^{i\theta}b$ and
$f_\sigma\to e^{i\theta}f_\sigma$ leave $X^{0\sigma}$ invariant:
the common phase is redundant. Discarding the constraint introduces,
among others, an auxiliary vacuum and a two-spinon state with no
counterparts in the projected atom. Enforcing it only on average is
a mean-field approximation. The idempotent $P_G$ and the Hubbard
operators above implement the projection exactly and offer local
tests for any auxiliary-particle calculation.

\subsection{Partial projection as a variational wave function}
What is lost if one discards doublons altogether? Return to the dimer,
with $t>0$ and exactly two particles. The free bonding determinant is
an equal-amplitude superposition of the singly occupied singlet
$\ket{s}$ and the symmetric doublon $\ket{D_+}$, in phases chosen
so that $\langle s|H|D_+\rangle=-2t$. Multiply each configuration
with a doublon by $g$, $0\le g\le1$. This is a small-system version
of a Gutzwiller variational state:
\[
 \ket{\psi(g)}=\frac{\ket{s}+g\ket{D_+}}{\sqrt{1+g^2}},
 \qquad E(g)=\frac{-4tg+Ug^2}{1+g^2}.
\]
There is no approximation in evaluating this expectation value.
More unusually, on the dimer this one-parameter family contains the
exact ground state. Stationarity gives $2tg^2+Ug-2t=0$, whose
positive root can be written without subtracting nearly equal numbers:
\begin{wl}
gopt = 4 t/(U + Sqrt[U^2 + 16 t^2]);
energy[y_] := (-4 t y + U y^2)/(1 + y^2);
FullSimplify[{energy[gopt] - emin,
  gopt^2/(2 (1 + gopt^2)) - dsite}, U > 0 && t > 0]
\end{wl}
The second check is the probability of finding a doublon on one
specified site. At $t/U\ll1$ the optimum has $g\simeq2t/U$, whereas
strict projection $g=0$ gives $E=0$ when evaluated with the original
Hamiltonian. A nonzero admixture of \emph{amplitude} $t/U$ produces
exchange energy of order $t^2/U$. The projected spin Hamiltonian
recovers that energy by transforming the Hamiltonian, rather than
by retaining the doublon in the state. For larger clusters inequivalent
hole--doublon separations generally require different amplitudes; one
scalar $g$ is no longer enough to make the trial state exact.

\section{A three-orbital atom: Hund's rule from the spectrum}
Distinguish a single-particle level index from a site index: now
$a_{m\sigma}$ removes an electron in orbital $m=1,2,3$ of \emph{one}
atom. Set all three orbital energies to zero. For a rotationally
invariant Kanamori interaction the repulsion in distinct orbitals is
$U'=U-2J$ for opposite spin and $U'-J$ for equal spin. Density terms
alone would give different energies to the three components of a spin
triplet. Two additional processes restore the multiplet structure:
\begin{align}
 H_K={}&U\sum_m D_m+
 (U-2J)\sum_{m\ne n} n_{m\uparrow}n_{n\downarrow}
 +(U-3J)\sum_{m<n,\sigma} n_{m\sigma}n_{n\sigma}\nonumber\\
 &-J\sum_{m\ne n} S_m^+ S_n^-
 +J\sum_{m\ne n} B_m^\dagger B_n,\qquad
 B_m=a_{m\downarrow}a_{m\uparrow}.
 \label{eq:kanamori}
\end{align}
The ordered last sums already contain their Hermitian conjugates. The
pair convention matters: reversing only the two annihilators would
reverse the pair-hopping amplitude. To make all factors explicit,
construct the terms for each unordered orbital pair:
\begin{wl}
snegfermionoperators[a];
snegrealconstants[J];
orb = a /@ Range[3];
n[m_, s_] := number[a[m], s];
pair[m_] := nc[a[AN, m, DO], a[AN, m, UP]];
hkpair[m_, k_] := (U - 2 J) (nc[n[m, UP], n[k, DO]] +
    nc[n[m, DO], n[k, UP]]) +
  (U - 3 J) Total[nc[n[m, #], n[k, #]] & /@ {UP, DO}] -
  J (nc[spinplus[a[m]], spinminus[a[k]]] +
     nc[spinminus[a[m]], spinplus[a[k]]]) +
  J (nc[conj[pair[m]], pair[k]] + nc[conj[pair[k]], pair[m]]);
hk = U Total[hubbard[orb]] + Total[hkpair @@@ Subsets[Range[3], {2}]];
Expand[{hk - conj[hk], komutator[hk, Total[spinplus /@ orb]]}]
\end{wl}

Put two electrons in six spin-orbitals. Here $Q=N-3=-1$. The
following rows give $S$, the energy (in symbolic $U,J$), and the
number of \emph{multiplets} at that energy. We keep the multiplicities;
distinct eigenvalues alone cannot diagnose a missing state.
\begin{wl}
bk = qsbasisvc[orb];
k2 = makematricesbzvc[hk, Select[bk, #[[1, 1]] == -1 &]];
levels = Flatten[Table[{b[[1, 2]], #[[1]], #[[2]]} & /@
    Tally[Expand /@ Eigenvalues[b[[2]]]], {b, k2}], 1];
TableForm[Sort[levels], TableHeadings -> {None, {"Spin", "Energy", "Multiplets"}}]
\end{wl}
For $J>0$ the lowest energy is $U-3J$. The two electrons occupy
different orbitals with parallel spins; three choices of orbital pair
give three triplets. Among singlets, each distinct-orbital pair gives
one state at $U-J$. The three intra-orbital pairs form a matrix with
diagonal $U$ and off-diagonal $J$. Its symmetric vector has energy
$U+2J$, and the two orthogonal vectors have energy $U-J$. Thus the
dimensions are $9+5+1=\binom62$. Dropping pair hopping would leave
three degenerate states at $U$ instead of this $1+2$ splitting.
\begin{wl}
Total[(2 #[[1]] + 1) #[[3]] & /@ levels] == Binomial[6, 2]
\end{wl}
\begin{furtherreading}
The atomic splittings determine which spin and orbital degrees of
freedom survive when the shell couples to a bath. For the three-orbital
impurity with this interaction see Horvat, Žitko and Mravlje~\cite{horvat}.
The \href{https://auger.ijs.si/nano/biblio/biblio.html}{authors' publication list}
provides access to the paper.
\end{furtherreading}

\section*{Exercises}
\addcontentsline{toc}{section}{Exercises}
\begin{exercise}[Counting charge defects]\label{ex:defects}
At $t=0$ count the half-filled ground states and the states exactly
$U$ higher on $L\ge2$ sites. How many states must a charge-only basis
contain at each energy for $L=4$?
\end{exercise}

\begin{exercise}[Detuned sites]\label{ex:detuned}
Add $\delta(n_2-n_1)/2$ to the half-filled dimer. Derive the exchange
$J(\delta)$ to second order in $t$ and state the condition on
$U,\delta,t$ for this spin-only description. Compare with a singlet
eigenvalue, taking $U=10$, $\delta=6$ and $t=1/10$ in common energy units.
\end{exercise}

\begin{exercise}[Addition versus spin gap]\label{ex:gaps}
At half filling find the dimer charge gap for all $U\ge0$, $t>0$.
Why does it remain positive at $U=0$? Compare it with the spin gap
at $U=0$ and at $U/t\gg1$.
\end{exercise}

\begin{exercise}[A spectral sum rule]\label{ex:moments}
Derive the first three high-frequency coefficients of the atomic
propagator at general filling. Which of them distinguish the two
Hubbard poles from a single Hartree pole? Verify with nested
commutators without extracting the poles a second time.
\end{exercise}
